% --- Archon physics formalization source begin ---
% archon:physics
% archon:covers PhyXMiniProblems/problem_phyx_mini_0878.lean
% archon:source-report reports/phyx_mini/problem_phyx_mini_0878.source.json

\chapter{Physics problem phyx\_mini\_0878}
\label{ch:PhyXMiniProblems_problem_phyx_mini_0878}

\paragraph{Problem source.}
\#\# Question

What is the equivalent capacitance of the three capacitors?

\#\# Answer choices

- A: A: 2.5\textbackslash{} \textbackslash{}mu\textbackslash{}mathrm\{F\}

- B: B: 35\textbackslash{} \textbackslash{}mu\textbackslash{}mathrm\{F\}

- C: C: 55\textbackslash{} \textbackslash{}mu\textbackslash{}mathrm\{F\}

- D: D: 25\textbackslash{} \textbackslash{}mu\textbackslash{}mathrm\{F\}

\#\# Figure caption (auxiliary; use the image as primary evidence)

The figure is a circuit diagram featuring three capacitors and a battery. Here's a detailed description:

1. **Battery:** Located on the left side, represented by a vertical line pair connected to the circuit, with positive (+) on the top and negative (−) on the bottom.

2. **Capacitors:**
   - Two capacitors, labeled as 20 μF and 30 μF, are arranged in parallel to each other and are connected across the battery.
   - Another capacitor, labeled as 13 μF, is connected in series with the parallel combination of the 20 μF and 30 μF capacitors.

3. **Circuit Layout:**
   - The battery is connected to the parallel combination of the 20 μF and 30 μF capacitors.
   - This parallel combination is then connected to the 13 μF capacitor, completing the circuit. 

The diagram displays a common configuration used in electric circuits to vary capacitance through series and parallel connections.

\#\# Dataset metadata

Electric Circuits; Physical Model Grounding Reasoning, Multi-Formula Reasoning

\paragraph{Recorded answer/context.}
D: D: 25\textbackslash{} \textbackslash{}mu\textbackslash{}mathrm\{F\}

\paragraph{Figure/image path.}
/root/proposal\_for\_physic/science-mango/phyx\_mini\_run/phyx\_data/test\_image/878.png

\paragraph{Formalization target.}
create a compiling Lean file with sorry bodies at `PhyXMiniProblems/problem\_phyx\_mini\_0878.lean`. The Lean declarations must preserve the physical quantities, dimensions or dimensional roles, figure labels, governing-law hypotheses, and final relation expressed by this problem.
Use Mathlib/Physlib names found through LeanExplore where available. If a physics API is missing, introduce faithful local abstractions rather than scalar placeholder aliases.

\begin{theorem}[Physics formalization target]
\label{thm:physics:phyx_mini_0878:target}
\lean{PhyXMiniProblems.ProblemPhyXMini0878.recorded_answer_is_D}
\uses{def:physics:phyx-mini-0878:phyxminiproblems-problemphyxmini0878-capacitanceinmicrofarads, def:physics:phyx-mini-0878:phyxminiproblems-problemphyxmini0878-equivalentcapacitancelaws, def:physics:phyx-mini-0878:phyxminiproblems-problemphyxmini0878-threecapacitorcircuit, def:physics:phyx-mini-0878:phyxminiproblems-problemphyxmini0878-matchesprimaryfigure, def:physics:phyx-mini-0878:phyxminiproblems-problemphyxmini0878-representsfigurenetworkreduction, def:physics:phyx-mini-0878:phyxminiproblems-problemphyxmini0878-answerchoice, def:physics:phyx-mini-0878:phyxminiproblems-problemphyxmini0878-displayedchoiceinmicrofarads}
The assigned autoformalize agent should translate this physics problem into Lean declarations in the covered file, with theorem and lemma proof bodies written as `by sorry`.
\end{theorem}
\begin{proof}
This is an autoformalization task, not a proof task. Produce statements that can later be proved by the physics prover without weakening the physical model.
\end{proof}

% --- Archon declaration topology begin ---
\section{Declaration topology}
\begin{definition}[capacitance Dimension]
  \label{def:physics:phyx-mini-0878:phyxminiproblems-problemphyxmini0878-capacitancedimension}
  \lean{PhyXMiniProblems.ProblemPhyXMini0878.capacitanceDimension}
  Capacitance has dimension C² T² M⁻¹ L⁻², since one farad is one coulomb per volt and voltage has dimension energy per charge
\end{definition}
\begin{proof}
  This is the declared model definition of capacitance Dimension.
\end{proof}
\begin{definition}[Capacitance Quantity]
  \label{def:physics:phyx-mini-0878:phyxminiproblems-problemphyxmini0878-capacitancequantity}
  \lean{PhyXMiniProblems.ProblemPhyXMini0878.CapacitanceQuantity}
  \uses{def:physics:phyx-mini-0878:phyxminiproblems-problemphyxmini0878-capacitancedimension}
  A nonnegative, unit-independent physical capacitance
\end{definition}
\begin{proof}
  This is the declared model definition of Capacitance Quantity.
\end{proof}
\begin{definition}[capacitance Readout]
  \label{def:physics:phyx-mini-0878:phyxminiproblems-problemphyxmini0878-capacitancereadout}
  \lean{PhyXMiniProblems.ProblemPhyXMini0878.capacitanceReadout}
  \uses{def:physics:phyx-mini-0878:phyxminiproblems-problemphyxmini0878-capacitancequantity}
  Read a physical capacitance in an arbitrary coherent system of units
\end{definition}
\begin{proof}
  This is the declared model definition of capacitance Readout.
\end{proof}
\begin{definition}[capacitance In Farads]
  \label{def:physics:phyx-mini-0878:phyxminiproblems-problemphyxmini0878-capacitanceinfarads}
  \lean{PhyXMiniProblems.ProblemPhyXMini0878.capacitanceInFarads}
  \uses{def:physics:phyx-mini-0878:phyxminiproblems-problemphyxmini0878-capacitancequantity, def:physics:phyx-mini-0878:phyxminiproblems-problemphyxmini0878-capacitancereadout}
  Read a physical capacitance in coherent-SI farads
\end{definition}
\begin{proof}
  This is the declared model definition of capacitance In Farads.
\end{proof}
\begin{definition}[capacitance In Microfarads]
  \label{def:physics:phyx-mini-0878:phyxminiproblems-problemphyxmini0878-capacitanceinmicrofarads}
  \lean{PhyXMiniProblems.ProblemPhyXMini0878.capacitanceInMicrofarads}
  \uses{def:physics:phyx-mini-0878:phyxminiproblems-problemphyxmini0878-capacitancequantity, def:physics:phyx-mini-0878:phyxminiproblems-problemphyxmini0878-capacitanceinfarads}
  Read a physical capacitance in microfarads
\end{definition}
\begin{proof}
  This is the declared model definition of capacitance In Microfarads.
\end{proof}
\begin{definition}[Circuit Node]
  \label{def:physics:phyx-mini-0878:phyxminiproblems-problemphyxmini0878-circuitnode}
  \lean{PhyXMiniProblems.ProblemPhyXMini0878.CircuitNode}
  The three electrically relevant nodes visible in image 878.png
\end{definition}
\begin{proof}
  This is the declared model definition of Circuit Node.
\end{proof}
\begin{definition}[Capacitor Label]
  \label{def:physics:phyx-mini-0878:phyxminiproblems-problemphyxmini0878-capacitorlabel}
  \lean{PhyXMiniProblems.ProblemPhyXMini0878.CapacitorLabel}
  Figure-derived names for the three physical capacitors
\end{definition}
\begin{proof}
  This is the declared model definition of Capacitor Label.
\end{proof}
\begin{definition}[Ideal Battery]
  \label{def:physics:phyx-mini-0878:phyxminiproblems-problemphyxmini0878-idealbattery}
  \lean{PhyXMiniProblems.ProblemPhyXMini0878.IdealBattery}
  \uses{def:physics:phyx-mini-0878:phyxminiproblems-problemphyxmini0878-circuitnode}
  The ideal battery symbol shown on the left of the image
\end{definition}
\begin{proof}
  This is the declared model definition of Ideal Battery.
\end{proof}
\begin{definition}[Capacitor]
  \label{def:physics:phyx-mini-0878:phyxminiproblems-problemphyxmini0878-capacitor}
  \lean{PhyXMiniProblems.ProblemPhyXMini0878.Capacitor}
  \uses{def:physics:phyx-mini-0878:phyxminiproblems-problemphyxmini0878-capacitancequantity, def:physics:phyx-mini-0878:phyxminiproblems-problemphyxmini0878-circuitnode}
  A physical capacitor together with its two endpoints. upperTerminal and lowerTerminal refer only to the orientation in the raster circuit diagram
\end{definition}
\begin{proof}
  This is the declared model definition of Capacitor.
\end{proof}
\begin{definition}[Series Capacitance Reduction]
  \label{def:physics:phyx-mini-0878:phyxminiproblems-problemphyxmini0878-seriescapacitancereduction}
  \lean{PhyXMiniProblems.ProblemPhyXMini0878.SeriesCapacitanceReduction}
  \uses{def:physics:phyx-mini-0878:phyxminiproblems-problemphyxmini0878-capacitancequantity}
  A claimed two-capacitance series reduction, before applying its law
\end{definition}
\begin{proof}
  This is the declared model definition of Series Capacitance Reduction.
\end{proof}
\begin{definition}[Parallel Capacitance Reduction]
  \label{def:physics:phyx-mini-0878:phyxminiproblems-problemphyxmini0878-parallelcapacitancereduction}
  \lean{PhyXMiniProblems.ProblemPhyXMini0878.ParallelCapacitanceReduction}
  \uses{def:physics:phyx-mini-0878:phyxminiproblems-problemphyxmini0878-capacitancequantity}
  A claimed two-branch parallel reduction, before applying its law
\end{definition}
\begin{proof}
  This is the declared model definition of Parallel Capacitance Reduction.
\end{proof}
\begin{definition}[Equivalent Capacitance Laws]
  \label{def:physics:phyx-mini-0878:phyxminiproblems-problemphyxmini0878-equivalentcapacitancelaws}
  \lean{PhyXMiniProblems.ProblemPhyXMini0878.EquivalentCapacitanceLaws}
  \uses{def:physics:phyx-mini-0878:phyxminiproblems-problemphyxmini0878-capacitancereadout, def:physics:phyx-mini-0878:phyxminiproblems-problemphyxmini0878-seriescapacitancereduction, def:physics:phyx-mini-0878:phyxminiproblems-problemphyxmini0878-parallelcapacitancereduction}
  The elementary physical laws for reducing ideal capacitors. They are stated at every coherent unit choice; within each equation all readouts therefore use one common coordinate and have the same dimensional role
\end{definition}
\begin{proof}
  This is the declared model definition of Equivalent Capacitance Laws.
\end{proof}
\begin{definition}[Three Capacitor Circuit]
  \label{def:physics:phyx-mini-0878:phyxminiproblems-problemphyxmini0878-threecapacitorcircuit}
  \lean{PhyXMiniProblems.ProblemPhyXMini0878.ThreeCapacitorCircuit}
  \uses{def:physics:phyx-mini-0878:phyxminiproblems-problemphyxmini0878-capacitorlabel, def:physics:phyx-mini-0878:phyxminiproblems-problemphyxmini0878-idealbattery, def:physics:phyx-mini-0878:phyxminiproblems-problemphyxmini0878-capacitor, def:physics:phyx-mini-0878:phyxminiproblems-problemphyxmini0878-seriescapacitancereduction, def:physics:phyx-mini-0878:phyxminiproblems-problemphyxmini0878-parallelcapacitancereduction}
  The circuit quantities to be reduced. The two equivalent capacitances here are unknown physical quantities, not numerical answers
\end{definition}
\begin{proof}
  This is the declared model definition of Three Capacitor Circuit.
\end{proof}
\begin{definition}[Matches Primary Figure]
  \label{def:physics:phyx-mini-0878:phyxminiproblems-problemphyxmini0878-matchesprimaryfigure}
  \lean{PhyXMiniProblems.ProblemPhyXMini0878.MatchesPrimaryFigure}
  \uses{def:physics:phyx-mini-0878:phyxminiproblems-problemphyxmini0878-capacitanceinmicrofarads, def:physics:phyx-mini-0878:phyxminiproblems-problemphyxmini0878-threecapacitorcircuit}
  Endpoint and printed-value observations read from the primary image. In particular, these fields encode the centre *series* branch shown in the image, not the reversed topology asserted by the auxiliary caption
\end{definition}
\begin{proof}
  This is the declared model definition of Matches Primary Figure.
\end{proof}
\begin{definition}[Represents Figure Network Reduction]
  \label{def:physics:phyx-mini-0878:phyxminiproblems-problemphyxmini0878-representsfigurenetworkreduction}
  \lean{PhyXMiniProblems.ProblemPhyXMini0878.RepresentsFigureNetworkReduction}
  \uses{def:physics:phyx-mini-0878:phyxminiproblems-problemphyxmini0878-threecapacitorcircuit}
  The reduction order inferred from the endpoint topology: first reduce the two centre capacitors in series, then put that result in parallel with the right capacitor. This predicate contains no numerical value for either output
\end{definition}
\begin{proof}
  This is the declared model definition of Represents Figure Network Reduction.
\end{proof}
\begin{definition}[Answer Choice]
  \label{def:physics:phyx-mini-0878:phyxminiproblems-problemphyxmini0878-answerchoice}
  \lean{PhyXMiniProblems.ProblemPhyXMini0878.AnswerChoice}
  Multiple-choice labels in the order printed in the problem
\end{definition}
\begin{proof}
  This is the declared model definition of Answer Choice.
\end{proof}
\begin{definition}[displayed Choice In Microfarads]
  \label{def:physics:phyx-mini-0878:phyxminiproblems-problemphyxmini0878-displayedchoiceinmicrofarads}
  \lean{PhyXMiniProblems.ProblemPhyXMini0878.displayedChoiceInMicrofarads}
  \uses{def:physics:phyx-mini-0878:phyxminiproblems-problemphyxmini0878-answerchoice}
  The capacitance value displayed beside each choice, in microfarads
\end{definition}
\begin{proof}
  This is the declared model definition of displayed Choice In Microfarads.
\end{proof}
\begin{lemma}[center series equivalent eq twelve microfarads]
  \label{lem:physics:phyx-mini-0878:phyxminiproblems-problemphyxmini0878-center-series-equivalent-eq-twelve-microfarads}
  \lean{PhyXMiniProblems.ProblemPhyXMini0878.center_series_equivalent_eq_twelve_microfarads}
  \uses{def:physics:phyx-mini-0878:phyxminiproblems-problemphyxmini0878-capacitanceinmicrofarads, def:physics:phyx-mini-0878:phyxminiproblems-problemphyxmini0878-equivalentcapacitancelaws, def:physics:phyx-mini-0878:phyxminiproblems-problemphyxmini0878-threecapacitorcircuit, def:physics:phyx-mini-0878:phyxminiproblems-problemphyxmini0878-matchesprimaryfigure, def:physics:phyx-mini-0878:phyxminiproblems-problemphyxmini0878-representsfigurenetworkreduction}
  The two centre capacitors reduce to 12 μF
\end{lemma}
\begin{proof}
  The conclusion follows from the governing relations and figure readouts named in the statement of center series equivalent eq twelve microfarads.
\end{proof}
\begin{theorem}[equivalent capacitance eq twenty five microfarads]
  \label{thm:physics:phyx-mini-0878:phyxminiproblems-problemphyxmini0878-equivalent-capacitance-eq-twenty-five-microfarads}
  \lean{PhyXMiniProblems.ProblemPhyXMini0878.equivalent_capacitance_eq_twenty_five_microfarads}
  \uses{def:physics:phyx-mini-0878:phyxminiproblems-problemphyxmini0878-capacitanceinmicrofarads, def:physics:phyx-mini-0878:phyxminiproblems-problemphyxmini0878-equivalentcapacitancelaws, def:physics:phyx-mini-0878:phyxminiproblems-problemphyxmini0878-threecapacitorcircuit, def:physics:phyx-mini-0878:phyxminiproblems-problemphyxmini0878-matchesprimaryfigure, def:physics:phyx-mini-0878:phyxminiproblems-problemphyxmini0878-representsfigurenetworkreduction}
  the equivalent capacitance of all three capacitors is 25 μF
\end{theorem}
\begin{proof}
  The conclusion follows from the governing relations and figure readouts named in the statement of equivalent capacitance eq twenty five microfarads.
\end{proof}
% --- Archon declaration topology end ---
% --- Archon physics formalization source end ---
